\section{Adaptive Gamification Engine}

\label{sec:gamification}
% ══════════════════════════════════════════════════════════════════════════════

\subsection{Challenge Generation}

The gamification engine generates personalized challenges for each
observer. A challenge $c$ specifies: target taxon group, geographic
area, time window, and observation protocol. Challenges are designed
to fill gaps in the existing data coverage.

\subsection{Reinforcement Learning Formulation}

We model challenge generation as a contextual bandit problem. The
\textbf{state} $s_t$ for observer $i$ at time $t$ comprises:
\begin{itemize}
  \item Observer profile: skill level, taxonomic expertise, location,
    historical observation patterns.
  \item Data gap map: spatial-temporal-taxonomic cells with low
    observation density.
  \item Environmental context: season, weather, habitat accessibility.
\end{itemize}

The \textbf{action} $a_t$ is a challenge parameterized by
$(\text{taxon}, \text{location}, \text{time}, \text{difficulty})$.

The \textbf{reward} $r_t$ combines:
\begin{equation}
  r_t = w_1 \cdot \text{DataValue}(a_t) + w_2 \cdot \text{Engagement}(a_t)
  + w_3 \cdot \text{Quality}(a_t)
  \label{eq:reward}
\end{equation}

where DataValue measures the scientific value of filling the targeted
gap, Engagement measures whether the observer attempted and completed
the challenge, and Quality measures the reliability of the resulting
observation.

\subsection{Optimization}

We use a contextual Thompson sampling algorithm
\citep{agrawal2013thompson} with a neural network reward model:
\begin{equation}
  \hat{r}(s, a; \theta) = \text{MLP}_\theta([\mathbf{s} \| \mathbf{a}])
\end{equation}

The reward model is updated daily on the previous day's challenge
outcomes. Posterior uncertainty is maintained through an ensemble of
5 models with different random initializations.

\subsection{Challenge Types}

\begin{table}[H]
\centering
\caption{Challenge types and their gamification mechanisms.}
\label{tab:challenges}
\begin{tabular}{p{2.5cm}p{4cm}p{4.5cm}}
\toprule
\textbf{Challenge} & \textbf{Objective} & \textbf{Gamification} \\
\midrule
Explorer & Visit undersampled 1\,km grid cell & Map reveal, territory
  badge \\
Species Quest & Find target species in area & Progressive difficulty,
  species album \\
Survey Sprint & Complete structured survey & Streak bonuses, team
  competitions \\
Night Watch & Nocturnal survey in target area & Exclusive night badge,
  bonus multiplier \\
Seasonal Tracker & Document phenological events & Season completion
  progress bar \\
Data Rescue & Verify/correct historical records & Accuracy leaderboard \\
\bottomrule
\end{tabular}
\end{table}

\subsection{Bias Reduction Results}

\begin{table}[H]
\centering
\caption{Sampling bias metrics: ZooPlatform vs.\ unguided observation.}
\label{tab:bias}
\begin{tabular}{lcc}
\toprule
\textbf{Bias metric} & \textbf{Unguided} & \textbf{ZooPlatform} \\
\midrule
Spatial Gini coefficient & 0.83 & 0.44 \\
Weekend/weekday ratio & 3.7:1 & 1.8:1 \\
Bird/invertebrate obs.\ ratio & 14.2:1 & 5.1:1 \\
Distance from road (median km) & 0.4 & 1.8 \\
\% of 1\,km cells with $\geq$1 obs. & 8.3\% & 22.7\% \\
\bottomrule
\end{tabular}
\end{table}

Adaptive gamification reduces spatial Gini coefficient by 47\% (0.83 to
0.44), taxonomic bias by 64\% (14.2:1 to 5.1:1), and increases spatial
coverage by 2.7x.


% ══════════════════════════════════════════════════════════════════════════════
